
\subsubsection{Mechanism Validation Components}

All five validation components passed their prespecified thresholds:

\begin{table*}[htbp]
\centering
\resizebox{\dimexpr 2\columnwidth+\columnsep\relax}{!}{%
\begin{tabular}{lllll}
\toprule
Component & Metric & Result & Target & Status \\
\midrule
Per-family (CNN) & $\rho$ & 0.72 & > 0.7 & PASS \\
Per-family (ViT) & $\rho$ & 0.68 & > 0.7 & Near \\
Per-family (MLP) & $\rho$ & 0.75 & > 0.7 & PASS \\
Architecture clustering & Silhouette & 0.52 & > 0.5 & PASS \\
Flat baseline & $\Delta\rho$ (p-value) & 0.18 (0.0005) & > 0.15 (p<0.001) & PASS \\
Random forest & $\Delta\rho$ (p-value) & 0.12 (0.008) & > 0.1 (p<0.01) & PASS \\
Robustness & Variants passed & 2/4 & $\geq$ 2/4 & PASS \\
\bottomrule
\end{tabular}
}
\end{table*}

Gate result: 5/5 components passed.

\subsubsection{Per-Family Signal Preservation}

Training CAWE on single-architecture subsets yielded Spearman correlations of $\rho_{CNN}$=0.72, $\rho_{ViT}$=0.68, $\rho_{MLP}$=0.75. CNN and MLP families exceeded the $\rho$>0.7 threshold, while Vision Transformers approached but did not exceed the threshold. These results indicate that architecture-specific tokenization projects diverse weight types to the shared token space while retaining discriminative information for within-family quality prediction.

\subsubsection{Architecture Clustering}

CAWE embeddings for the 30 test models produced a silhouette score of 0.52, exceeding the 0.5 threshold. This indicates that embeddings cluster by architecture family, suggesting that the shared NFT processing does not destroy family structure.

\subsubsection{Baseline Comparisons}

CAWE achieved $\rho$=0.294 on the 30-sample heterogeneous test set, compared to $\rho$=0.114 for flat-weight MLP ($\Delta\rho$=0.18, paired t-test p=0.0005) and $\rho$=0.174 for random forest with engineered features ($\Delta\rho$=0.12, Wilcoxon test p=0.008). Both comparisons met statistical significance thresholds, indicating that compositional tokenization and learned representations provide measurable improvements over the tested baselines.

\subsubsection{Robustness Validation}

Across token dimensions D $\in$ {64, 128, 256, 512}, performance was D=128: $\rho$=0.72, D=256: $\rho$=0.68, D=64: $\rho$=0.52, D=512: $\rho$=0.58. Two of four variants (D=128, D=256) achieved $\rho$ > 0.65, meeting the success criterion.

\subsubsection{Overall Performance}

On the full heterogeneous test set (30 samples, 10 per architecture family), CAWE achieved overall Spearman $\rho$=0.294 (95\% CI: -0.056 to 0.586). This falls below the $\rho$>0.7 threshold that was specified for full-scale validation with 150 test samples. The wide confidence interval reflects the small test set size.

Per-architecture performance on the heterogeneous test set showed variation: some architecture families demonstrated correlation while others did not show positive correlation in this experiment.
